% Half-title, title page and colophon.
\providecommand{\bookversion}{}
\newcommand{\bookdate}{\number\day\space\ifcase\month\or January\or February\or March\or April\or
  May\or June\or July\or August\or September\or October\or November\or December\fi\space
  \number\year}
\newcommand{\bookedition}{\ifx\bookversion\empty\else\bookversion\,·\,\fi\bookdate}

\thispagestyle{empty}
\vspace*{2.2in}
\begin{flushright}
  {\Large\bfseries CAN}
\end{flushright}
\cleardoublepage

\thispagestyle{empty}
\vspace*{1.1in}
\begin{flushleft}
  {\color{accent}\rule{0.9in}{2.5pt}}\par\vspace{1.4ex}
  {\fontsize{32}{38}\selectfont\bfseries CAN\par}
  \vspace{2.2ex}
  {\Large\itshape The bus, the frame and the controller\par}
  \vspace{1.2in}
  {\large Erik Pihl\par}
  \vfill
  \vspace{0.6ex}
  {\small\color{muted} \bookedition\par}
\end{flushleft}
\newpage

\thispagestyle{empty}
\vspace*{\fill}
{\small
\noindent\textit{CAN - the bus, the frame and the controller}\\
A short book about CAN, written for whoever is about to build a CAN node: either the controller
in hardware or the driver in software.

\medskip
\noindent The book describes CAN as a protocol: the bus, the frame, bit stuffing, the checksum,
arbitration and bit timing. The last chapter describes the simplified controller that is built and
driven in teaching, what it implements of all that comes before, and what it deliberately leaves
out.

\medskip
\noindent This is the English edition of
\foreignlanguage{swedish}{\textit{CAN - bussen, framen och kontrollern}}. The text is in English,
and so are all the code, all the register names and all the signal names, exactly as in the code
and in the shared specifications.

\medskip
\noindent Text and figures: Erik Pihl.

\medskip
\noindent The book's text and figures, like the course material it is typeset from, are licensed
under Creative Commons Attribution 4.0 International (CC BY 4.0). You may copy, distribute and
adapt the material for any purpose, commercial ones included, provided that you credit the author
and state the licence. The licence is at \url{https://creativecommons.org/licenses/by/4.0/}.

\medskip
\noindent Typeset with Lua\LaTeX\ in TeX Gyre Pagella and DejaVu Sans Mono.

\medskip
\noindent\textcolor{muted}{This edition: \bookedition.}
\par}
\cleardoublepage
